% ============================================================
%  ICN-292 · Sistemas de Información para la Gestión · 2026-2
%  Entrega 1 (E1) — Propuesta SIG para una PYME
%  Caso: IMELSE
%
%  Compilar con: pdflatex (dos veces, por el índice y las refs)
%  Recomendado: Overleaf
%
%  ESTRUCTURA OBLIGATORIA (secciones A–G del enunciado).
%  No cambiar el orden ni los títulos: el corrector revisa por ellos.
%
%============================================================

\documentclass[10pt,letterpaper]{article}

\usepackage[utf8]{inputenc}
\usepackage[T1]{fontenc}
\usepackage[spanish,es-noquoting]{babel}   % en Overleaf funciona; si falla localmente, comentar
\usepackage[margin=2cm]{geometry}
\usepackage{graphicx}
\usepackage{booktabs}
\usepackage{longtable}
\usepackage{array}
\usepackage{float}
\usepackage{pdflscape}      % páginas apaisadas para los BPMN
\usepackage{caption}
\usepackage{enumitem}
\usepackage{fancyhdr}
\usepackage{lastpage}
\usepackage[hidelinks]{hyperref}

\pagestyle{fancy}
\fancyhf{}
\lhead{\footnotesize ICN-292 · Entrega 1 · IMELSE}
\rhead{\footnotesize \thepage\ / \pageref{LastPage}}
\renewcommand{\headrulewidth}{0.4pt}

\captionsetup{font=small,labelfont=bf}
\setlist{itemsep=2pt,parsep=0pt}

% ---------- Datos editables ----------
\newcommand{\paralelo}{P100}
\newcommand{\repo}{https://github.com/IMicit/ICN292\_T3\_20262}
% -------------------------------------

\begin{document}

% ============================================================
% A. PORTADA Y FICHA DEL EQUIPO
% ============================================================
\begin{titlepage}
\centering
\vspace*{2cm}

{\Large\bfseries Propuesta de Sistema de Información para la Gestión}\\[0.4cm]
{\Large Trazabilidad de mantención preventiva de tableros eléctricos}\\[0.3cm]
{\large IMELSE}\\[1.5cm]

\rule{\textwidth}{0.8pt}\\[0.4cm]
{\large Entrega 1 (E1)}\\[0.2cm]
{\normalsize ICN-292 Sistemas de Información para la Gestión · Paralelo 100 · Segundo Semestre 2026}\\[0.2cm]
\rule{\textwidth}{0.8pt}\\[1.5cm]

\begin{tabular}{@{}llll@{}}
\toprule
\textbf{Nombre} & \textbf{RUT} & \textbf{ROL USM} & \textbf{Rol en el equipo} \\
\midrule
Javier Bravo    & 21.330.540-9 & 202360638-3 & Modelo de datos \\
Benjamín Eliz   & 21.803.009-2 & 202360632-4 & Revisión y autenticación \\
Cristóbal Muñoz & 21.657.535-0 & 202360603-0 & Documentación y QA \\
Nicolás Oyarzún & 21.762.466-5 & 202404599-7 & Registro en terreno \\
Nicolás Soto & 22.083.026-8 & 202460585-2 & Panel e informes \\
\bottomrule
\end{tabular}

\vspace{1.2cm}
\textbf{Repositorio GitHub:}\\
\url{\repo}

\vfill
{\small Universidad Técnica Federico Santa María · Departamento de Industrias · Campus Vitacura\\
Profesores: José Miguel González Paul · José Luis Sáez Tamayo\\
Fecha de entrega: 12 de septiembre de 2026}
\end{titlepage}

\tableofcontents
\newpage

% ============================================================
% B. PYME Y PROBLEMA (EVIDENCIA)
% ============================================================
\section{PYME y problema}

\subsection{Identificación de la PYME}

\begin{table}[H]
\centering
\small
\begin{tabular}{@{}p{4.5cm}p{10cm}@{}}
\toprule
\textbf{Campo} & \textbf{Dato} \\
\midrule
Nombre de fantasía & IMELSE \\
Razón social / RUT & 76.046.641-7 \\
Rubro & Servicios de mantenimiento eléctrico industrial, foco retail \\
Constitución & Persona natural desde 2000; sociedad de responsabilidad limitada desde 2009 \\
Ubicación & Av. General Saavedra 1217, Independencia, Región Metropolitana \\
Contacto & (56-2) 2732 3122 · contacto@imelse.cl · \url{https://imelse.cl} \\
Gerente General & Yuri Rojas \\
Dotación & 24 personas: 14 fijas y 10 estudiantes en formación dual (convenio SOFOFA) \\
Distribución & 4 en oficina, 20 en terreno \\
Clientes principales & SMU (Unimarc, bajo contrato), Cencosud (Jumbo, Santa Isabel), Falabella (Tottus, Sodimac) \\
Cobertura geográfica & Región Metropolitana: Tiltil por el norte, Linderos por el sur, Talagante por el poniente \\
\bottomrule
\end{tabular}
\caption{Identificación de la empresa}
\end{table}

\paragraph{Justificación del alcance PYME.}
IMELSE cuenta con una dotación de 24 personas, de las cuales 14 corresponden a
personal fijo y 10 a estudiantes en formación dual, y opera con una estructura
administrativa de cuatro personas en oficina. La toma de decisiones operativas se
concentra en la Gerencia General y en la Gerencia de Operaciones, sin capas
intermedias de supervisión ni área de tecnologías de información. La empresa no
dispone de sistemas de gestión propios: la coordinación se apoya en planillas,
formularios impresos y mensajería instantánea. Este perfil corresponde al de una
pequeña empresa según el criterio de tamaño de la Ley 20.416, y define además la
restricción de diseño principal de esta propuesta: cualquier solución debe poder ser
operada y mantenida sin personal técnico dedicado.

\subsection{Evidencia de existencia}

\begin{table}[H]
\centering
\small
\begin{tabular}{@{}clll@{}}
\toprule
\textbf{N°} & \textbf{Evidencia} & \textbf{Fecha} & \textbf{Referencia} \\
\midrule
1 & Sitio web institucional \url{https://imelse.cl} & 01/09/2026 & \texttt{assets/ev-01-sitio.png} \\
2 & Acta de entrevista firmada por el Gerente de Operaciones & 01/09/2026 & \texttt{assets/ev-02-acta.pdf} \\
3 & Fichas de checklist y hoja de servicio en uso (anonimizadas) & 06/09/2026 & \texttt{assets/ev-05-ficha-checklist.pdf} \\
4 & Registro fotográfico de la visita a dependencias & 06/09/2026 & \texttt{assets/ev-06-visita.pdf} \\
5 & Transcripción completa de la entrevista & 01/09/2026 & \texttt{docs/03-transcripcion-entrevista.pdf} \\
\bottomrule
\end{tabular}
\caption{Evidencia verificable del caso}
\end{table}

\paragraph{Entrevista.}
La caracterización del proceso y del problema proviene de una entrevista
semiestructurada al Gerente de Operaciones de IMELSE, Francisco Meza, realizada el
01/09/2026 en dependencias de la empresa. La pauta utilizada, el acta firmada y la transcripción
completa se encuentran en el repositorio. El entrevistado autorizó el uso académico
del material en el marco de este curso.

\paragraph{Nota de confidencialidad.}
Los documentos operativos facilitados por la empresa se incorporan anonimizados:
se ocultaron nombres y RUT de trabajadores, códigos y direcciones de locales del
cliente corporativo y referencias a condiciones comerciales.

\subsection{Problema de negocio}

\begin{quote}
\itshape
Entre 15 y 20 de cada 100 mantenciones preventivas quedan incompletas sin detección
oportuna, porque el cumplimiento se declara en checklists de papel sin evidencia
verificable por tablero y llega a oficina 16 a 24 horas después. Esto genera
re-visitas absorbidas por la empresa, tableros que pasan ciclos bimensuales completos
sin revisión, y exposición a falla eléctrica bajo un contrato con seguro comprometido
de 27.500 UF.
\end{quote}

\paragraph{Nota metodológica.} La cifra de 15 a 20 por ciento corresponde a una
estimación declarada por el Gerente de Operaciones durante la entrevista y no a una
medición. Se cita como tal, y su conversión en indicador medido es precisamente uno
de los resultados esperados del sistema propuesto.

\paragraph{Manifestaciones.}
\begin{enumerate}[label=P\arabic*.,leftmargin=*]
  \item \textbf{Traspaso de información diferido.} El resultado de la visita llega a
  oficina entre 16 y 24 horas después de ejecutada, cuando la cuadrilla regresa a
  dependencias de IMELSE y entrega los formularios en papel.
  \item \textbf{Registro no verificable.} El cumplimiento se declara marcando un
  checklist impreso, sin evidencia asociada a cada tablero, por lo que no es posible
  distinguir un tablero intervenido de uno solo declarado.
  \item \textbf{Cobertura de tableros incierta.} Existe levantamiento de tableros en
  el 80 a 90 por ciento de los locales, pero no se contrasta contra lo efectivamente
  ejecutado, de modo que un tablero puede atravesar ciclos bimensuales completos sin
  revisión sin que nadie lo detecte.
  \item \textbf{Pérdida de evidencia.} Las fotografías de respaldo se envían por
  mensajería instantánea y posteriormente no se encuentran o no pueden asociarse a la
  visita y al tablero que documentan.
  \item \textbf{Detección tardía y reactiva.} Los incumplimientos se descubren en la
  revisión posterior del informe o ante una falla en el local, no durante la
  ejecución.
  \item \textbf{Re-trabajo en abastecimiento.} La información incompleta sobre el
  estado de los tableros genera compras erradas y retiros de materiales que no
  corresponden a la necesidad real.
\end{enumerate}

\paragraph{Impacto.}

\begin{table}[H]
\centering
\small
\begin{tabular}{@{}p{2.6cm}p{6.2cm}p{5.6cm}@{}}
\toprule
\textbf{Dimensión} & \textbf{Manifestación} & \textbf{Consecuencia para IMELSE} \\
\midrule
Servicio & Mantenciones preventivas incompletas (P2, P3) & Incumplimiento del alcance comprometido en el contrato bimensual \\
Costo & Re-visitas para completar trabajo declarado (P2, P5) & Costo de cuadrilla y traslado absorbido íntegramente por la empresa \\
Tiempo & Latencia de 16 a 24 horas en el traspaso (P1) & La gerencia decide sobre información del día anterior \\
Riesgo & Tableros sin revisión durante ciclos completos (P3) & Exposición a falla eléctrica bajo un seguro comprometido de 27.500 UF \\
Control & Ausencia de evidencia asociada al tablero (P2, P4) & Imposibilidad de auditar el cumplimiento ante el cliente corporativo \\
Abastecimiento & Información incompleta del estado real (P6) & Compras erradas y retiros de bodega innecesarios \\
\bottomrule
\end{tabular}
\caption{Impacto del problema por dimensión}
\end{table}

\paragraph{Intentos previos de solución.}
La empresa ensayó dos mecanismos antes de este proyecto, y ambos fallaron por la
misma razón estructural.

El primero consistió en instalar una etiqueta adhesiva en cada tablero intervenido,
como marca física de ejecución. El mecanismo se abandonó porque los técnicos dejaban
de instalarla ante cualquier faltante o condición irregular, para no comprometer
responsabilidad sobre un trabajo que consideraban incompleto. El registro dependía de
la voluntad de quien lo generaba y, en el momento en que más importaba, desaparecía.

El segundo, vigente desde hace aproximadamente un mes y medio, es una hoja de registro
de tableros con mediciones de termografía. Mejora la calidad del diagnóstico técnico,
pero conserva las dos limitaciones centrales: sigue siendo un documento en papel que
llega a oficina con la misma latencia, y no garantiza que la hoja cubra todos los
tableros del local.

Ambos intentos muestran que el problema no es de diligencia del personal sino de
diseño del registro: mientras la declaración de cumplimiento se pida después del
hecho y sin evidencia asociada, el mecanismo es evitable. Esto justifica que la
solución propuesta deba \textbf{forzar el registro en el momento de la ejecución} en
lugar de solicitarlo posteriormente, decisión que atraviesa todo el rediseño de
proceso y el modelo de datos de este informe.

\subsection{Objetivo del SIG propuesto}

\begin{quote}
\itshape
Permitir que la gerencia determine, en el momento del cierre de cada visita y sin
depender de la declaración del técnico, qué tableros del local fueron efectivamente
intervenidos y con qué resultado, sustituyendo el checklist de papel por un registro
guiado con evidencia obligatoria por tablero.
\end{quote}

El objetivo se acota deliberadamente a una decisión concreta y verificable. No se
propone un sistema de gestión integral: se propone cerrar la brecha entre lo
declarado y lo ejecutado en la mantención preventiva bajo contrato, que es donde se
concentran el riesgo contractual y el costo de re-visita.

% ============================================================
% C. REQUERIMIENTOS DESDE CERO
% ============================================================
\section{Requerimientos}

\subsection{Actores y roles}

\begin{table}[H]
\centering
\small
\begin{tabular}{@{}p{3.6cm}p{2.2cm}p{8.8cm}@{}}
\toprule
\textbf{Actor} & \textbf{Tipo} & \textbf{Relación con el sistema} \\
\midrule
Técnico de terreno & Humano & Ejecuta la mantención y registra cada tablero con evidencia \\
Ayudante & Humano & Apoya la ejecución; puede registrar bajo el mismo perfil \\
Estudiante dual & Humano & Integra la cuadrilla en formación; perfil de registro restringido \\
Programación / oficina & Humano & Confirma el ciclo, asigna cuadrilla y fecha \\
Administración y bodega & Humano & Gestiona la entrega de materiales y el stock \\
Gerencia de Operaciones & Humano & Consulta el panel, revisa cobertura, pendientes y hallazgos \\
Cliente corporativo & Externo & Recibe el aviso de servicio y el informe de la visita \\
Encargado del local & Externo & Valida la presencia de la cuadrilla mediante firma en el dispositivo \\
Plataforma SIG & Sistema & Genera visitas, valida cobertura, cierra visitas y genera informes \\
\bottomrule
\end{tabular}
\caption{Actores y roles}
\end{table}

\subsection{Alcance}

\subsubsection*{Dentro del alcance}
\begin{itemize}
  \item Mantención preventiva bajo contrato, en la rama del cliente corporativo SMU.
  \item Catálogo maestro de clientes, locales y tableros.
  \item Programación de visitas a partir del ciclo bimensual del contrato.
  \item Registro guiado por tablero con evidencia fotográfica obligatoria.
  \item Validación de cobertura y declaración de pendientes con motivo.
  \item Panel de gerencia y generación del informe de visita.
\end{itemize}

\subsubsection*{Fuera del alcance}

\begin{table}[H]
\centering
\small
\begin{tabular}{@{}p{5cm}p{9.5cm}@{}}
\toprule
\textbf{Exclusión} & \textbf{Justificación} \\
\midrule
Servicios correctivos & Se originan por aviso y no tienen alcance predefinido, por lo que la validación de cobertura no aplica. Se modelan como contexto en el as-is \\
Proyectos y mejoras de infraestructura & Operan con lógica de cotización y no de ciclo de contrato \\
Gestión de inventario y bodega & Existe un procedimiento en uso que no forma parte del problema declarado \\
Facturación y cobranza & Fuera del objetivo del sistema; no interviene en la brecha entre lo declarado y lo ejecutado \\
Integración con sistemas del cliente & Depende de terceros y de acuerdos comerciales no disponibles en el horizonte del proyecto \\
Equipos de frío & El alcance técnico de IMELSE llega hasta el punto eléctrico \\
\bottomrule
\end{tabular}
\caption{Exclusiones de alcance y su justificación}
\end{table}

\subsection{Requisitos funcionales}

\begin{longtable}{@{}p{1.4cm}p{8.4cm}p{1.8cm}p{1.6cm}p{1.4cm}@{}}
\toprule
\textbf{ID} & \textbf{Requisito} & \textbf{Actor} & \textbf{Prioridad} & \textbf{Origen} \\
\midrule
\endhead
RF-01 & El sistema debe mantener un registro maestro de clientes, locales y tableros, con identificación única de cada tablero dentro de su local & Programación & M & P3 \\
RF-02 & El sistema debe generar la visita a partir del calendario del contrato y asociarle una cuadrilla y una fecha & Programación & M & P1 \\
RF-03 & El sistema debe presentar un checklist guiado por tablero que no permita marcar un ítem como cumplido sin la evidencia asociada & Técnico & M & P2 \\
RF-04 & El sistema debe exigir al menos una fotografía por tablero intervenido, capturada durante la ejecución & Técnico & M & P4 \\
RF-05 & El sistema debe registrar automáticamente fecha, hora y usuario en cada registro, sin intervención del operador & Plataforma & M & P2 \\
RF-06 & El sistema debe impedir el cierre de la visita si existen tableros del local sin registro, salvo declaración explícita de pendiente & Plataforma & M & P3 \\
RF-07 & El sistema debe permitir declarar un tablero como pendiente indicando obligatoriamente el motivo & Técnico & M & P3 \\
RF-08 & El sistema debe dejar el resultado de la visita disponible en oficina de forma inmediata al cierre & Plataforma & M & P1 \\
RF-09 & El sistema debe mantener el historial de intervenciones por tablero a lo largo de los ciclos & Gerencia & M & P5 \\
RF-10 & El sistema debe registrar los hallazgos detectados durante la visita, clasificados por nivel de criticidad & Técnico & S & P5 \\
RF-11 & El sistema debe registrar los materiales o repuestos faltantes detectados en terreno, asociados al local y a la visita & Técnico & S & P6 \\
RF-12 & El sistema debe entregar a la gerencia un panel con el porcentaje de cobertura de tableros y de visitas completas por local y por ciclo & Gerente & S & P3, P5 \\
RF-13 & El sistema debe generar el informe de mantención en formato exportable, con la evidencia asociada, para su envío al cliente & Administración & S & P4 \\
RF-14 & El sistema debe administrar usuarios y perfiles de acceso diferenciados para técnico, gerencia y administración & Gerente & S & --- \\
RF-15 & El sistema debe arrastrar automáticamente los pendientes declarados de un ciclo a la planificación del ciclo siguiente & Gerente & C & P3 \\
RF-16 & El sistema debe permitir al encargado del local validar la presencia de la cuadrilla mediante firma en el dispositivo & Encargado & C & --- \\
RF-17 & El sistema debe registrar las mediciones numéricas de termografía asociadas a cada tablero & Técnico & C & P5 \\
RF-18 & El sistema debe notificar automáticamente a la gerencia cuando una visita se cierra con pendientes críticos & Plataforma & C & P5 \\
\bottomrule
\caption{Requisitos funcionales priorizados (MoSCoW), trazados a las manifestaciones
del problema. El detalle de cada requisito se encuentra en
\texttt{docs/01-requerimientos.md}.}
\end{longtable}

\subsection{Requisitos no funcionales}

\begin{longtable}{@{}p{1.6cm}p{8.2cm}p{2.4cm}p{1.6cm}@{}}
\toprule
\textbf{ID} & \textbf{Requisito} & \textbf{Categoría} & \textbf{Prioridad} \\
\midrule
\endhead
RNF-01 & El registro debe poder completarse sin conexión y sincronizarse al recuperar señal, dado que las salas de tableros pueden carecer de cobertura celular & Disponibilidad & M \\
RNF-02 & Un registro cerrado no puede modificarse; las correcciones generan un registro nuevo con autor y fecha, conservando el original & Integridad & M \\
RNF-03 & El registro de un tablero debe completarse en un máximo de cinco interacciones, siendo operable con una mano y con guantes de seguridad & Usabilidad & M \\
RNF-04 & El sistema debe soportar el volumen de evidencia del ciclo: aproximadamente 30 visitas mensuales con unos 20 tableros cada una, con al menos una fotografía por tablero & Rendimiento & M \\
RNF-05 & El acceso debe estar controlado por perfil, y cada acción sobre un registro debe quedar atribuida a un usuario identificado & Seguridad & M \\
RNF-06 & El tratamiento de datos personales de trabajadores y de terceros debe ajustarse a la Ley 21.719, con minimización de datos y finalidad declarada & Cumplimiento & M \\
RNF-07 & El sistema debe ejecutarse en dispositivos móviles Android de gama media, sin requerir hardware especializado & Compatibilidad & S \\
RNF-08 & Las fotografías deben comprimirse antes de la sincronización, sin perder legibilidad de la identificación del tablero ni de las conexiones & Rendimiento & S \\
RNF-09 & El prototipo de la Entrega 2 debe ser reproducible en localhost a partir del repositorio, con instrucciones en el README & Mantenibilidad & M \\
RNF-10 & El sistema debe conservar la evidencia por al menos el período de vigencia del contrato de mantención & Cumplimiento & S \\
\bottomrule
\caption{Requisitos no funcionales priorizados (MoSCoW). El detalle de cada requisito
se encuentra en \texttt{docs/01-requerimientos.md}.}
\end{longtable}

Dos precisiones. La exigencia de operación sin conexión (RNF-01) depende de la
disponibilidad de señal celular en las salas eléctricas, consulta enviada a la empresa
y sin respuesta al cierre de este informe; se mantiene como \emph{Must} por criterio
conservador y se declara fuera del alcance del prototipo en la sección
\ref{sec:alcance-prototipo}. El límite de cinco interacciones (RNF-03) responde a que
cerca del 40 por ciento del personal de apoyo son estudiantes duales en rotación, por
lo que el sistema no puede suponer entrenamiento acumulado.

\subsection{Trazabilidad problema $\rightarrow$ requisito}

\begin{table}[H]
\centering
\small
\begin{tabular}{@{}p{1.2cm}p{6.3cm}p{6.8cm}@{}}
\toprule
\textbf{ID} & \textbf{Manifestación} & \textbf{Requisitos que la atacan} \\
\midrule
P1 & Traspaso de información diferido & RF-02, RF-08 \\
P2 & Registro no verificable & RF-03, RF-04, RF-05, RNF-02 \\
P3 & Cobertura de tableros incierta & RF-01, RF-06, RF-07, RF-12, RF-15 \\
P4 & Pérdida de evidencia & RF-04, RF-05, RF-13 \\
P5 & Detección tardía y reactiva & RF-06, RF-08, RF-09, RF-10, RF-18 \\
P6 & Re-trabajo en abastecimiento & RF-09, RF-11 \\
\bottomrule
\end{tabular}
\caption{Trazabilidad entre las manifestaciones del problema y los requisitos}
\end{table}

% ============================================================
% D. PROCESOS BPMN
% ============================================================
\section{Modelado de procesos}

\subsection{Alcance del modelo}

El modelo as-is incluye la recepción de órdenes de trabajo y la compuerta de tipo de
trabajo con sus ramas de servicios correctivos y de proyectos, con el fin de situar el
proceso dentro de la operación completa de IMELSE. El modelado en profundidad y el
rediseño se acotan a la rama de \textbf{mantención preventiva bajo contrato},
coherente con el alcance declarado en la sección anterior. Las ramas de correctivos y
proyectos se representan sin desarrollo interno y no son objeto de rediseño.

Los archivos fuente en formato Bizagi (\texttt{bpmn-as-is-imelse.bpm} y
\texttt{bpmn-to-be-imelse.bpm}) están disponibles en \texttt{assets/} del repositorio,
junto con las imágenes exportadas.

\subsection{Proceso as-is}

\paragraph{Puntos críticos identificados.}
El tramo de cierre, representado en la Figura \ref{fig:asis}, concentra las
debilidades del proceso actual. La generación del informe está condicionada al
retorno físico de la cuadrilla a dependencias de IMELSE, y un evento temporizador de
16 a 24 horas separa la ejecución en terreno de la disponibilidad de la información en
oficina. El proceso contiene además dos ciclos de re-trabajo: el de solicitud de
pruebas fotográficas, que puede obligar a la cuadrilla a volver al local a tomar
nuevas fotografías, y el de solicitud de detalles faltantes tras la revisión del
informe por parte de Gerencia.

Un tercer punto crítico no es un elemento del diagrama sino su ausencia: no existe
ninguna compuerta que verifique que los tableros intervenidos correspondan a los
tableros existentes en el local. El cumplimiento se declara en el formulario
\emph{Lista de chequeo} sin contraste posible, y los pendientes se comunican
verbalmente, sin quedar registrados ni corregidos, tal como consigna la anotación
respectiva en la Figura \ref{fig:asis}.

\clearpage
\begin{figure}[p]
  \centering
  \vspace*{\fill}
  \includegraphics[angle=90, width=0.92\textwidth, height=0.88\textheight, keepaspectratio]{bpmn-as-is-imelse.png}
  \caption{Proceso as-is: desde la recepción de la orden hasta el cierre e informe}
  \label{fig:asis}
  \vspace*{\fill}
\end{figure}
\clearpage

\subsection{Proceso to-be}

\paragraph{Nota de notación.} El subproceso \emph{Realizar mantenimiento por tablero}
se representa con marcador de múltiple instancia secuencial, indicando que se ejecuta
una vez por cada tablero del local. La tarea \emph{Consultar panel}, en el carril de
Gerencia, no recibe flujo de secuencia: se activa por un evento temporizador propio y
se conecta al almacén mediante asociación de datos. Esta representación es
deliberada y expresa que la revisión de gerencia dejó de ser un paso bloqueante del
proceso operativo. Las cuatro representaciones del almacén \emph{Sistema de registro
SIG} corresponden al mismo repositorio lógico, y se repiten únicamente por legibilidad
del diagrama.

\begin{figure}[H]
  \centering
  \includegraphics[angle=90, width=0.48\textheight, height=0.95\textwidth, keepaspectratio]{bpmn-to-be-imelse.png}
  \caption{Proceso to-be con la intervención del SIG}
  \label{fig:tobe}
\end{figure}

\clearpage

\subsection{Explicación de las mejoras}

\paragraph{Qué se automatiza.} La generación de la visita desde el calendario del
contrato, el aviso de servicio al cliente, la marca de fecha, hora y usuario de cada
tablero intervenido, la generación del informe y su envío al cliente. En el proceso
actual todas estas dependen del retorno físico de la cuadrilla a las oficinas.

\paragraph{Qué se controla.} La cobertura de tableros, mediante una compuerta previa
al cierre que compara el levantamiento del local contra los tableros efectivamente
registrados en la visita. La visita no puede cerrarse con tableros sin intervenir
salvo declaración explícita de pendiente con motivo. El control se ejerce en el
momento y no en la revisión posterior, que es precisamente lo que fallaron los dos
intentos previos de la empresa.

\paragraph{Qué se mide.} Tres variables que hoy no existen como dato: cobertura de
tableros por visita, tasa de mantención completa y latencia entre la ejecución en
terreno y la disponibilidad del registro en oficina.

\paragraph{Qué desaparece.} El evento temporizador de 16 a 24 horas, los dos ciclos de
re-trabajo asociados a la solicitud de fotografías y de detalles faltantes, la
dependencia del retorno a dependencias de IMELSE para generar el informe, y la salida
informal por pendiente sin registro, que pasa a ser un estado registrado con motivo y
autor.

\begin{table}[H]
\centering
\begin{tabular}{lcc}
\toprule
\textbf{Elemento del modelo} & \textbf{as-is} & \textbf{to-be} \\
\midrule
Tareas & 31 & 20 + 1 subproceso \\
Compuertas (decisiones / uniones) & 9 (5 / 4) & 6 (3 / 3) \\
Eventos intermedios & 6 & 2 \\
\quad de ellos, eventos de espera & 3 & 2 \\
Eventos de término & 1 & 2 \\
Bucles de re-trabajo & 2 & 0 \\
Carriles con actividad & 3 & 4 \\
Tareas en el carril del sistema & 0 & 7 \\
Puntos de control de cobertura & 0 & 1 \\
Asociaciones de datos & 4 & 7 \\
Latencia del registro & 16--24 h & inmediata \\
\bottomrule
\end{tabular}
\caption{Comparación estructural entre el as-is y el to-be, sobre la ruta de
mantención preventiva bajo contrato.}
\label{tab:comparacion-asis-tobe}
\end{table}

Dos precisiones sobre el cuadro. Los dos eventos de espera que subsisten están ambos
en el bloque de solicitud de materiales, que no fue rediseñado; los eventos de espera
eliminados corresponden en su totalidad al tramo de cierre. Los eventos de término
aumentan de uno a dos porque la revisión de gerencia dejó de ser un paso bloqueante
del flujo y pasó a constituir un fragmento independiente, conectado al registro por
asociación de datos.

% ============================================================
% E. MODELO DE DATOS PRELIMINAR
% ============================================================
\clearpage
\section{Modelo de datos preliminar}

\begin{figure}[H]
  \centering
  \includegraphics[angle=90, width=0.88\textwidth, height=0.78\textheight, keepaspectratio]{er-preliminar.png}
  \caption{Modelo entidad-relación preliminar}
  \label{fig:er}
\end{figure}

\clearpage

\subsection{Entidades y claves}

\begin{table}[H]
\centering
\small
\begin{tabular}{@{}p{2.4cm}p{2.8cm}p{3.2cm}p{5.6cm}@{}}
\toprule
\textbf{Entidad} & \textbf{PK} & \textbf{FK} & \textbf{Atributos principales} \\
\midrule
Cliente & id\_cliente & --- & razon\_social, unidad\_negocio \\
Local & id\_local & id\_cliente & nombre, dirección, comuna \\
Tablero & id\_tablero & id\_local & código, ubicación, tipo \\
Ciclo & id\_ciclo & --- & año, periodo, fecha\_inicio \\
Visita & id\_visita & id\_local, id\_ciclo & fecha, estado \\
Usuario & id\_usuario & --- & nombre, rol \\
Cuadrilla & id\_visita + id\_usuario & --- & función \\
ItemChecklist & id\_item & --- & descripción, orden, obligatorio \\
RegistroItem & id\_registro & id\_visita, id\_tablero, id\_item, id\_usuario & resultado, ts\_captura \\
Evidencia & id\_evidencia & id\_registro & archivo, ts, hash \\
Pendiente & id\_pendiente & id\_visita, id\_tablero & motivo, comentario \\
Hallazgo & id\_hallazgo & id\_registro & criticidad, descripción \\
\bottomrule
\end{tabular}
\caption{Entidades del modelo preliminar y sus claves}
\end{table}

\paragraph{Relaciones.} Todas las relaciones del modelo son de tipo uno a muchos, con
la excepción de \texttt{Cuadrilla}, entidad asociativa que resuelve la relación de
muchos a muchos entre \texttt{Visita} y \texttt{Usuario} y que permite registrar
además la función desempeñada por cada integrante en esa visita específica.

\paragraph{La relación que sostiene la solución.} La relación \texttt{Local}
uno a muchos \texttt{Tablero} es la que hace posible el requisito central del sistema.
Sin un catálogo de los tableros que pertenecen a cada local no existe conjunto contra
el cual contrastar lo ejecutado, y la validación de cobertura de RF-06 carece de
referencia. Es también la razón por la cual la construcción del catálogo maestro
(RF-01) es condición previa de toda la propuesta y no un requisito accesorio.

\paragraph{Clave primaria de \texttt{Tablero}.} Se evaluaron dos alternativas: una
clave natural compuesta por \texttt{id\_local} y el código rotulado en el tablero, y
una clave surrogate \texttt{id\_tablero}. Se optó por la surrogate, con una
restricción de unicidad sobre el par \texttt{(id\_local, código)} que preserva la
regla de negocio de que no existen dos tableros con el mismo código dentro de un
mismo local.

La justificación es doble. En primer lugar, el requisito no funcional de
inmutabilidad (RNF-02) establece que un registro cerrado no puede modificarse; si la
clave primaria fuera el código físico, cualquier reetiquetado, reemplazo de tablero o
corrección del rotulado obligaría a propagar el cambio sobre registros históricos ya
cerrados, contradiciendo dicho requisito. La clave surrogate permite que el código
evolucione como atributo sin afectar la trazabilidad acumulada por tablero, que es lo
que exige RF-09.

En segundo lugar, el levantamiento de tableros existe actualmente en el 80 a 90 por
ciento de los locales, de modo que parte del catálogo debe poder crearse antes de
disponer de un código definitivo y verificado en terreno. Una clave natural exigiría
un código válido en el momento del alta.

Como consideración adicional, la clave natural es compuesta y se propagaría como par
de columnas hacia \texttt{RegistroItem} y \texttt{Pendiente}, ampliando las claves
foráneas de las entidades de mayor volumen del modelo sin beneficio funcional, dado
que la validación de cobertura opera por identificador interno.

\subsection{Trazabilidad proceso $\leftrightarrow$ datos}

\begin{table}[H]
\centering
\small
\begin{tabular}{@{}p{6cm}p{8.5cm}@{}}
\toprule
\textbf{Tarea del to-be} & \textbf{Entidades que lee o escribe} \\
\midrule
Generar ciclo para la visita & lee \texttt{Local}, \texttt{Ciclo}; escribe \texttt{Visita} \\
Asignar cuadrilla & escribe \texttt{Cuadrilla} \\
Registrar tablero (subproceso) & lee \texttt{Tablero}, \texttt{ItemChecklist}; escribe \texttt{RegistroItem}, \texttt{Evidencia} \\
Registrar hallazgo & escribe \texttt{Hallazgo} \\
Validar cobertura de tableros & compara \texttt{Tablero} del local contra \texttt{RegistroItem} de la visita \\
Declarar pendiente & escribe \texttt{Pendiente} \\
Cerrar visita & actualiza el estado de \texttt{Visita} \\
Generar informe & lee \texttt{Visita}, \texttt{RegistroItem}, \texttt{Evidencia}, \texttt{Hallazgo} \\
Consultar panel & lee \texttt{RegistroItem}, \texttt{Pendiente}, \texttt{Hallazgo} \\
\bottomrule
\end{tabular}
\caption{Cómo el modelo de datos sostiene el proceso modelado}
\end{table}

\subsection{Decisiones y supuestos del modelo}

\paragraph{La entidad \emph{Informe} se dejó fuera del modelo.} El informe de visita
es un artefacto generado íntegramente a partir de \texttt{Visita},
\texttt{RegistroItem}, \texttt{Evidencia} y \texttt{Hallazgo}, y no aporta información
que no esté ya contenida en esas entidades. Incorporarlo como entidad introduciría
duplicación y un riesgo de divergencia entre el informe almacenado y los datos que lo
originaron. Se trata de una decisión de diseño y no de una omisión.

\paragraph{Las fotografías se almacenan como archivos.} La tabla \texttt{Evidencia}
conserva la ruta del archivo junto con su hash y la marca de tiempo de captura. El
hash permite detectar alteraciones posteriores del archivo, lo que complementa a
RNF-02: no basta con que el registro sea inmutable si la evidencia asociada puede
sustituirse.

\paragraph{La cuadrilla se modela como entidad asociativa.} Una visita involucra
varios usuarios y un usuario participa en múltiples visitas. La entidad
\texttt{Cuadrilla} resuelve esta relación y permite además registrar la función
específica de cada integrante, distinción necesaria porque el perfil de estudiante
dual tiene atribuciones de registro distintas a las del técnico responsable.

\paragraph{El checklist se modela como catálogo.} \texttt{ItemChecklist} es una
entidad independiente y no un conjunto de columnas de \texttt{RegistroItem}, de modo
que la pauta de mantención puede evolucionar sin alterar el esquema ni invalidar los
registros históricos.

% ============================================================
% F. ARQUITECTURA LÓGICA Y STACK
% ============================================================
\section{Arquitectura lógica y stack tentativo}

\subsection{Componentes lógicos}

La solución se organiza en cuatro capas. La \textbf{captura} se realiza mediante una
aplicación web responsiva usada desde el teléfono por la cuadrilla, con fotografía
tomada desde la cámara del dispositivo. La \textbf{persistencia} combina una base de
datos relacional con el esquema del modelo ER y almacenamiento de archivos para las
imágenes, referenciadas por ruta y hash. La capa \textbf{analítica} expone los
indicadores mediante vistas y consultas sobre cobertura de tableros, visitas completas
y latencia. La \textbf{interfaz} entrega el panel de gerencia y la generación del
informe de visita.

\subsection{Stack tentativo para la Entrega 2}

Se propone Python 3.12 con Django 5 y PostgreSQL 16, con SQLite como alternativa para
simplificar la reproducción por parte del corrector. Las imágenes se almacenan en el
sistema de archivos local y se referencian desde \texttt{Evidencia}, evitando
dependencias de servicios externos y credenciales. La interfaz se construye con
plantillas Django de diseño responsivo, de modo que un único desarrollo sirva al panel
de oficina y al uso en terreno.

El criterio de selección fue el dominio del equipo y la reproducibilidad, no la
superioridad técnica: dado que el prototipo debe levantarse en localhost por un
tercero y que el plazo es acotado, se privilegiaron herramientas conocidas y de
configuración mínima por sobre alternativas de mayor escalabilidad, que el volumen del
caso no requiere.

\subsection{Reproducibilidad en localhost}

El \texttt{README.md} documentará la puesta en marcha completa: clonar el repositorio,
crear el entorno virtual, instalar las dependencias de \texttt{requirements.txt},
aplicar las migraciones, cargar los datos ficticios de prueba y levantar el servidor.
Los datos de prueba incluirán un cliente, dos locales con su catálogo de tableros, un
ciclo y usuarios para cada rol, de modo que el flujo completo pueda recorrerse sin
carga manual previa.

\subsection{Alcance del prototipo}
\label{sec:alcance-prototipo}

El prototipo de la Entrega 2 implementará la captura en línea, la validación de
cobertura, la declaración de pendientes, el cierre de visita y el panel de gerencia
con los indicadores definidos.

Se declara explícitamente fuera del alcance del prototipo la operación sin conexión
(RNF-01). El requisito se contempla en el diseño, y el modelo de datos lo admite en
tanto el registro es una entidad independiente con marca de tiempo de captura propia,
pero su implementación exige una capa de sincronización y resolución de conflictos que
excede el plazo disponible. Se prefiere declarar esta restricción antes que
comprometerla y no cumplirla.

\begin{table}[H]
\centering
\small
\begin{tabular}{@{}p{5.5cm}p{3.5cm}p{3.5cm}@{}}
\toprule
\textbf{Indicador} & \textbf{Línea base} & \textbf{Meta} \\
\midrule
Tasa de mantención completa & 80--85 \% (declarado) & $\geq$ 95 \% \\
Cobertura de tableros por visita & No medible actualmente & $\geq$ 98 \% \\
Latencia del registro & 16--24 h & $<$ 1 h \\
Tasa de re-visita & No medido actualmente & $-$50 \% \\
\bottomrule
\end{tabular}
\caption{Indicadores a implementar en la Entrega 2}
\end{table}

\paragraph{Indicador descartado.} Se evaluó y descartó un indicador de productividad
individual por técnico. La razón es metodológica y ética a la vez: el problema
declarado es de diseño del registro y no de desempeño de las personas, como lo
demuestra el fracaso de los dos intentos previos, y un indicador por persona
convertiría una herramienta de control de proceso en un instrumento de evaluación
laboral. Esta decisión se retoma en la sección de riesgos de datos personales.

% ============================================================
% G. PLAN DE ENTREGA 2 Y RIESGOS
% ============================================================
\section{Plan de Entrega 2 y riesgos}

\subsection{Hitos y responsables}

\begin{table}[H]
\centering
\small
\begin{tabular}{@{}clp{5.5cm}p{3cm}@{}}
\toprule
\textbf{N°} & \textbf{Fecha} & \textbf{Hito} & \textbf{Responsable} \\
\midrule
1 & Semana 1 & Esquema de base de datos implementado y migrado & Javier Bravo \\
2 & Semana 1 & Carga del catálogo maestro y datos de prueba & Javier Bravo \\
3 & Semana 2 & Autenticación y perfiles por rol & Benjamín Eliz \\
4 & Semana 2 & Generación de visitas desde el ciclo & Benjamín Eliz \\
5 & Semana 3 & Registro guiado por tablero con evidencia obligatoria & Nicolás Oyarzún \\
6 & Semana 3 & Validación de cobertura y declaración de pendientes & Nicolás Oyarzún \\
7 & Semana 4 & Panel de gerencia e indicadores & Nicolás Soto \\
8 & Semana 4 & Generación del informe de visita & Nicolás Soto \\
9 & Semana 5 & Pruebas de reproducción en localhost por un integrante ajeno al desarrollo & Cristóbal Muñoz \\
10 & Semana 5 & Documentación final, README y entrega & Cristóbal Muñoz \\
\bottomrule
\end{tabular}
\caption{Hitos de la Entrega 2}
\end{table}

\subsection{Riesgos del proyecto}

\begin{table}[H]
\centering
\small
\begin{tabular}{@{}p{5.5cm}p{2cm}p{6.5cm}@{}}
\toprule
\textbf{Riesgo} & \textbf{Impacto} & \textbf{Mitigación} \\
\midrule
Ausencia de levantamiento de tableros en el 10 a 20 \% de los locales & Alto & Modelar el catálogo como incremental: el sistema permite dar de alta tableros durante la visita, registrando quién los incorporó \\
Falta de confirmación sobre cobertura celular en salas de tableros & Alto & Declarar la operación sin conexión fuera del alcance del prototipo y diseñar el registro con marca de tiempo de captura propia, que la admite posteriormente \\
Resistencia del personal de terreno al registro obligatorio & Medio & Limitar el registro a cinco interacciones por tablero (RNF-03) y comunicar que los indicadores son de proceso y no de persona \\
Indisponibilidad de la empresa para validar el prototipo & Medio & Validar contra la transcripción y el acta ya obtenidas; registrar los supuestos no confirmados como tales \\
Concentración del desarrollo en pocos integrantes & Medio & Asignación explícita de hitos y commits identificables por integrante en el repositorio \\
Alcance del prototipo mayor al plazo disponible & Alto & Implementar por prioridad MoSCoW: los nueve requisitos \emph{Must} antes que cualquier \emph{Should} \\
\bottomrule
\end{tabular}
\caption{Riesgos identificados}
\end{table}

\subsection{Riesgos de datos personales (Ley 21.719)}

La Ley 21.719, publicada en el Diario Oficial el 13 de diciembre de 2024, reforma el
marco de la Ley 19.628 y crea la Agencia de Protección de Datos Personales como
autoridad fiscalizadora. Su entrada en plena vigencia está fijada para el 1 de
diciembre de 2026, dentro del horizonte de operación del sistema propuesto, por lo que
el diseño debe considerarla desde ahora. Cabe consignar que a agosto de 2026 el
Gobierno manifestó estar evaluando una eventual postergación de esa fecha, asociada a
la conformación del Consejo Directivo de la Agencia; mientras no exista modificación
legal, el 1 de diciembre de 2026 sigue siendo el plazo vigente. Un aspecto relevante
para este caso es que la ley cubre los datos de trabajadores y no solamente los de
clientes.

\paragraph{Datos personales tratados.} Nombre y rol de los trabajadores de IMELSE que
integran las cuadrillas, incluidos los estudiantes en formación dual; la asociación
entre cada trabajador y los registros que genera, con marca de fecha y hora; y las
fotografías de evidencia, en las que puede aparecer incidentalmente personal del local
o del cliente.

\paragraph{Finalidad declarada.} Acreditar la ejecución de la mantención preventiva
por tablero ante el cliente corporativo y permitir a la gerencia de IMELSE verificar
la cobertura del contrato. La finalidad es de control de proceso y de respaldo
contractual.

\paragraph{Base de licitud.} El tratamiento se funda en la ejecución del contrato de
servicios con el cliente corporativo y en la relación laboral, en tanto el registro
documenta una obligación contractual de la empresa y la actividad se realiza en el
marco de las funciones del trabajador. No se recurre al consentimiento como base,
precisamente porque en una relación laboral el consentimiento difícilmente puede
considerarse libre.

\paragraph{Minimización.} El diseño incorpora tres decisiones de minimización. No se
captura la geolocalización del dispositivo, dado que la identificación del tablero ya
establece el lugar de ejecución del trabajo y la ubicación adicional no aporta a la
finalidad declarada. Las fotografías se acotan al equipo intervenido, con instrucción
explícita de no registrar personas. Y no se registran tiempos de permanencia más allá
de las marcas de captura estrictamente necesarias para calcular el indicador de
latencia.

\paragraph{Conservación.} Los registros y su evidencia se conservan mientras dure la
vigencia del contrato de servicios más el período de eventuales reclamaciones
asociadas, dado que su finalidad es acreditar el cumplimiento ante el cliente. Cumplido
ese plazo, los datos que identifican al trabajador pueden disociarse conservando el
registro técnico del tablero, que es lo que sostiene el historial de RF-09.

\paragraph{Riesgo específico.} El riesgo principal del sistema es que, por diseño,
registra qué trabajador ejecutó qué intervención y en qué momento. Aunque la finalidad
declarada sea el control del proceso, el mismo dato permite construir una evaluación
individual de desempeño, uso distinto del declarado. Las medidas adoptadas son cuatro:
transparencia previa hacia los trabajadores sobre qué se registra y para qué; acceso
restringido por rol, de modo que la vista individual no esté disponible en el panel de
gerencia; definición de indicadores de proceso y no de persona, decisión ya
materializada en el descarte del indicador de productividad individual; y la
recomendación de realizar una evaluación de impacto en protección de datos antes de la
puesta en producción, considerando que se trata de un tratamiento sistemático de datos
de trabajadores.

\paragraph{Medidas en el prototipo.} El prototipo de la Entrega 2 operará
exclusivamente con datos ficticios: no se cargarán nombres reales de trabajadores de
IMELSE, ni códigos o direcciones reales de locales del cliente corporativo, ni
fotografías tomadas en instalaciones reales. El control de acceso por rol se
implementará desde el inicio y no como agregado posterior, de modo que la restricción
de visibilidad sea verificable en la demostración.

% ============================================================
% ANEXOS
% ============================================================
\vspace{1.2cm}
\appendix
\section{Anexos}

\subsection{Evidencia fotográfica}

Las evidencias citadas en la Tabla 2 se encuentran en assets/ y docs/ del repositorio, en los archivos \texttt{ev-01-sitio.png}, \texttt{ev-02-acta.pdf},
\texttt{ev-05-ficha-checklist.pdf} y \texttt{ev-06-visita.pdf}.
\subsection{Estructura del repositorio}
\begin{verbatim}
README.md
docs/     00-caso-pyme.md  01-requerimientos.md  02-bpmn.md
          03-transcripcion-entrevista.pdf
assets/   bpmn-as-is-imelse.bpm   bpmn-as-is-imelse.png
          bpmn-to-be-imelse.bpm   bpmn-to-be-imelse.png
          er-preliminar.dbml      er-preliminar.png
          ev-01 a ev-06 (evidencia)
informe/  ICN292_P100_E1_Bravo_Eliz_Muñoz_Oyarzún_Soto.pdf
          ICN292_P100_E1_Bravo_Eliz_Muñoz_Oyarzún_Soto.tex
\end{verbatim}

\end{document}
